% GENERATED FILE. Edit canonical lessons, then run npm run generate:books.
% canonical-source-hash: fnv1a64:77993917c3517139
% canonical-lessons: HI-C12-pitaa-maataa, HI-C12-bhaai-bahin

\chapter{Family}
\label{ch:hi-family}

\begin{chapteropening}
\textbf{By the end of this chapter:} \emph{I can name my father, mother, brother and sister in Hindi.}

Say \hi{भाई} and \hi{बहन} beside \hi{पिता} and \hi{माता}, and name which one is a true cousin of English 'brother' and which is not.
\end{chapteropening}

\section[\hi{पिता} \hi{माता}]{\hi{पिता}, \hi{माता} (pitā, mātā) — father and mother, cousins across a huge family}

\label{lesson:HI-C12-pitaa-maataa}



\begin{quote}
You've already met Latin \textbf{pater/māter} and their English cousins \textbf{father/mother} — all descending from one ancient shared word. Hindi keeps a cousin of its own, in the very same family.
\end{quote}

\begin{cousinweb}
\begin{itemize}
\raggedright
  \item \textbf{\hi{पिता}} (\emph{pitā}) = \textquotedblleft{}\textbf{father}\textquotedblright{} — Sanskrit \textbf{\hi{पितृ}} (\emph{pitṛ}), from the \textbf{same PIE root}, \emph{\textbf{ph\textsubscript{2}t\'{\={e}}r}}, as Latin \emph{pater} and English \emph{father}.
  \item \textbf{\hi{माता}} (\emph{mātā}) = \textquotedblleft{}\textbf{mother}\textquotedblright{} — Sanskrit \textbf{\hi{मातृ}} (\emph{mātṛ}), from \emph{\textbf{méh\textsubscript{2}tēr}}, the same root as Latin \emph{māter} and English \emph{mother}.
\end{itemize}

So \emph{pitā} and \emph{father} are, at the deepest level, the \textbf{same ancient word} — reaching you here through Sanskrit rather than through Latin or Germanic, one more branch of the same very old family tree.
\end{cousinweb}

\subsection*{You'll want to know — Be honest about register — the everyday words}
\textbf{\hi{पिता}/\hi{माता}} are correct, respectable Hindi — but the words you'll hear \textbf{most often in daily life} are different: \textbf{\hi{बाप}} (\emph{bāp}, \textquotedblleft{}father,\textquotedblright{} an everyday/colloquial word) and \textbf{\hi{माँ}} (\emph{māṁ}, \textquotedblleft{}mother,\textquotedblright{} universally the warm, everyday word — closer to English \textquotedblleft{}mom\textquotedblright{} than \textquotedblleft{}mother\textquotedblright{}). These aren't uncertain outliers, either — they're \textbf{also} Sanskrit-descended, just by a different path than \emph{pitā/mātā}: \textbf{\hi{बाप}} traces through Prakrit to Sanskrit \textbf{\hi{वप्तृ}} (\emph{vaptṛ}, \textquotedblleft{}father,\textquotedblright{} a different Sanskrit root from \emph{pitṛ}), while \textbf{\hi{माँ}} is a genuine \textbf{doublet} of \emph{mātā} — both come from the very same Sanskrit \emph{mātṛ}, just worn down along a more colloquial Prakrit line.

\subsection*{Your turn}
Say these aloud:

\begin{itemize}
\raggedright
  \item \textquotedblleft{}pitā, mātā\textquotedblright{} — the Sanskrit-descended, formal pair
  \item \textquotedblleft{}bāp, māṁ\textquotedblright{} — the everyday, spoken pair
  \item the deep cousin — \textquotedblleft{}pitā and father are the same ancient PIE word\textquotedblright{}
\end{itemize}

\begin{tcolorbox}[breakable,colback=teal!4,colframe=teal!35!black,title={Before you move on}]
What are \hi{पिता} and \hi{माता} cousins of, in Latin and English? (\textbf{Pater/ father} and \textbf{māter/mother} — all from the same PIE roots.) What are the everyday, spoken words instead? (\textbf{\hi{बाप} bāp} and \textbf{\hi{माँ} māṁ}.) Are \emph{bāp} and \emph{māṁ} also Sanskrit-descended? (\textbf{Yes} — \emph{bāp} traces to Sanskrit \emph{vaptṛ}, a different root from \emph{pitṛ}; \emph{māṁ} is a genuine \textbf{doublet} of \emph{mātā}, both from \emph{mātṛ}.)
\end{tcolorbox}
\section[\hi{भाई} \hi{बहन}]{\hi{भाई}, \hi{बहन} (bhāī, bahin) — one cousin word, one different word entirely}

\label{lesson:HI-C12-bhaai-bahin}



\begin{quote}
You've met \emph{frāter} and \emph{soror} as PIE cousins of \textquotedblleft{}brother\textquotedblright{} and \textquotedblleft{}sister.\textquotedblright{} Hindi's pair has a genuine surprise hiding in it: only \textbf{one} of these two continues that ancient family.
\end{quote}

\subsection*{You'll want to know — \hi{भाई} — the expected cousin}
\textbf{\hi{भाई}} (\emph{bhāī}) = \textquotedblleft{}\textbf{brother}\textquotedblright{} — from Sanskrit \textbf{\hi{भ्रातृ}} (\emph{bhrātṛ}), which \textbf{does} continue the same PIE root, \emph{\textbf{b\textsuperscript{h}réh\textsubscript{2}tēr}}, as Latin \emph{frāter} and English \emph{brother}. This one is exactly what you'd expect: one more cousin in the same old family.

\subsection*{You'll want to know — \hi{बहन} — NOT the expected cousin}
Sanskrit did have a direct PIE-cognate word for \textquotedblleft{}sister\textquotedblright{} too: \textbf{\hi{स्वसृ}} (\emph{svasṛ}), the very cousin of Latin \emph{soror} and English \emph{sister}. But that word \textbf{faded from use}. In its place, Sanskrit built a \textbf{different} word, \textbf{\hi{भगिनी}} (\emph{bhaginī}) — and \emph{that} word, not \emph{svasṛ}, is the confirmed ancestor of Hindi's everyday \textbf{\hi{बहन}} (\emph{bahin}). (Where \emph{bhaginī} itself comes from is the less certain part: it's often explained as \textbf{\hi{भग}} (\emph{bhaga}, \textquotedblleft{}a share, good fortune\textquotedblright{}) plus a feminine suffix, \textquotedblleft{}the one who shares\textquotedblright{} — but Sanskritists don't all agree on this, and some suspect it's a reshaping of an older form avoiding an unrelated, indelicate sense of \emph{bhaga}. That deeper question is worth a linguist's extra check; the \emph{bahin $\neq$ svasṛ} finding itself is solid.)

So Hindi's brother/sister pair is \textbf{lopsided} in its ancestry: \emph{bhāī} genuinely continues the ancient PIE word; \emph{bahin} comes from an entirely different Sanskrit formation that \textbf{replaced} the PIE-cognate word before Hindi ever inherited it. (Sanskrit-to-Hindi sound history has plenty of non-obvious turns.)

\subsection*{Your turn}
Say these aloud:

\begin{itemize}
\raggedright
  \item \textquotedblleft{}bhāī\textquotedblright{} — brother — the true PIE cousin of \textquotedblleft{}brother\textquotedblright{}
  \item \textquotedblleft{}bahin\textquotedblright{} — sister — from bhaginī, \textquotedblleft{}the one who shares,\textquotedblright{} NOT from the PIE word for sister
  \item the asymmetry — one cousin, one completely different word
\end{itemize}

\begin{tcolorbox}[breakable,colback=teal!4,colframe=teal!35!black,title={Before you move on}]
Does \textbf{\hi{भाई}} continue the same ancient root as English \textquotedblleft{}brother\textquotedblright{}? (\textbf{Yes} — Sanskrit \emph{bhrātṛ} < PIE \emph{b\textsuperscript{h}réh\textsubscript{2}tēr}, same as Latin \emph{frāter}.) Does \textbf{\hi{बहन}} continue the same ancient root as English \textquotedblleft{}sister\textquotedblright{}? (\textbf{No} — Sanskrit's PIE-cognate word, \emph{svasṛ}, faded; \emph{bahin} comes from \textbf{\hi{भगिनी}} \emph{bhaginī}, an unrelated Sanskrit formation — often glossed \textquotedblleft{}the one who shares,\textquotedblright{} though \emph{bhaginī}'s own deeper origin is itself debated.)
\end{tcolorbox}
